% !TEX root = ../../tree_notes.tex
% =====================================================================
%  Chapter 13 opener -- "Chapter at a Glance", three reading paths + the
%  ladder staircase, THE running example (the same 10-node tree T as
%  Ch.12, edges 1-2,1-3,2-4,2-5,3-6,3-7,5-8,7-9,9-10) re-drawn with the
%  in/out DP state overlay (fig:c10-runtree-dp), the "who owns what"
%  comparison (DP cookbook here vs Ch.1 binary recursions vs Ch.2
%  bottom-up vs Ch.4 rooting vs Ch.12 metrics vs Ch.14 knapsack budgets
%  vs Ch.15 rerooting vs Ch.16 LCA), roadmap, recognition guide,
%  constraint->tool table, notation.
%  Exports fig:c10-runtree-dp and fig:c10-taxonomy.
%  All DP tables below are hand-computed on T; every later section
%  re-uses exactly these numbers.
% =====================================================================

\begin{chaptertheorybox}{Chapter at a Glance}
Chapters~1--3 put structure on \emph{values}: the BST orders them, the heap
keeps one extreme alive, the trie reads them symbol by symbol. Chapter~12
measured the tree's \emph{shape} -- depth, height, diameter, centroid, radius.
This chapter asks the third question: given a fixed root, \textbf{what does
every node owe you if you let it \emph{compute} from its children's answers?}
Four DP shapes, all built by the same two ingredients (a DFS with a parent
parameter, and a per-node \emph{state} that folds children into a
tuple):

\term{state per node} (\secref{c10-subtrees}): the tuple $\mathcal S(v)$
$\bigO{\deg(v)}$-folds its children's tuples through one rule -- the
general-tree sibling of \thmref{c01-height-recursion}; LC~333's
subtree-BST record lives here.
\term{in/out two-state} (\secref{c10-twostate}): keep two numbers per node
(take it or leave it) and combine over children -- maximum independent set,
vertex cover, weighted colourings; CF~1529C's Parsa endpoints are exactly
this shape.
\term{harvest at the peak} (\secref{c10-harvest}): every unordered pair $(u,v)$
has a unique topmost node $\ell = \mathrm{LCA}(u,v)$ and the path splits into
two arms meeting there; \propref{c10-pairtop} turns $\sum_{u<v} F(u,v)$ into
$\sum_{\ell} \text{combine over pairs of arms at }\ell$ -- this is
\thmref{c09-dpjoin} (Ch.12's diameter identity) generalised to \emph{any}
additive path functional.
\term{distance-indexed DP} (\secref{c10-distance}): state $= $ a histogram over
depths inside the subtree; the merge is a knapsack over children (budget
analysis is Chapter~14's, previewed only). CF~161D's distance-$k$ census and
this chapter's \exref{c10-k2} run here.

\textbf{You will learn to}
\begin{itemize}[nosep]
  \item state the bottom-up DP template on general trees and read the tuple
        off a problem statement (what to carry, why it closes under
        children); \hfill\lv{2}--\lv{3}
  \item recognise the in/out shape and drive it end-to-end on CF~1529C
        (weighted endpoint choice) and LC~333 (subtree-BST tuple); \hfill\lv{3}--\lv{4}
  \item state and prove \propref{c10-pairtop}: every path harvests at its
        peak, so counting paths reduces to combining arms at each $\ell$; \hfill\lv{4}
  \item hold the $k$-indexed histogram per node, merge children in
        $\bigO{k}$ each, and cold-write CF~161D's distance-$k$ count; \hfill\lv{4}--\lv{5}
  \item combine everything on the Summit CF~1076E: subtree-of-$v$, depth-cap
        at $d$, difference array on the tour -- the harvest-at-peak idea
        wearing an Euler skin. \hfill\lv{5}
\end{itemize}

\textbf{Prerequisites.} Chapter~12 (\chref{c09}) for the parent-parameter DFS
and the "join arms at the peak" fact \thmref{c09-dpjoin} this chapter
generalises; Chapter~4 (\chref{c04}) for rooting; the bottom-up recursion
pattern of Chapter~2. Rerooting, per-pair queries and knapsack budget
analysis are \emph{deliberately not assumed} -- they are Chapters~15, 16 and
14 respectively.
\end{chaptertheorybox}

% ---------------- Reading paths + ladder position ----------------
\begin{chaptertheorybox}{How to Read This Chapter (three paths)}
Level badges $\lv1$--$\lv6$: $\lv2$ is a one-line recurrence you write from
the definition, $\lv{3}$ is a solid LC-Medium coded straight from a section's
skeleton, $\lv{4}$ needs \emph{two} ideas composed (the in/out pair inside
the child-fold, or the peak-sum inside a histogram), $\lv{5}$ is the Summit,
where the state must survive a depth cap and a difference array at the same
time.

\textbf{Cookbook path ($\lv2$--$\lv{3}$, ``I want the DP shapes'').} Read
\secref{c10-subtrees} and \secref{c10-twostate}; do the solved LC~333 and
CF~1529C and their practices. You leave able to cold-write the child-fold
template and the in/out recurrence -- which covers the entire
"tree-DP warm-up" band and every house-robber-style variant you will be
asked in an interview.

\textbf{Path-pattern path ($\lv4$).} Start at \secref{c10-harvest}:
\propref{c10-pairtop} states the pair/peak decomposition that Ch.12's
\thmref{c09-dpjoin} proved once for diameter; carry the machinery into
\secref{c10-distance} and do the \exref{c10-k2} census before coding
CF~161D. This is the reading that discharges the three forward promises
Ch.16 has been holding.

\textbf{Full path (contest track).} Everything in order, then the mixed set,
then the Summit CF~1076E -- at which point the forward ropes are obvious:
budgeted child merges (Chapter~14, \chref{c27}), rerooting as theory
(Chapter~15, \chref{c20}), and per-pair distance \emph{queries}
(Chapter~16, \chref{c11}).

The backward ropes: Chapter~12 (\chref{c09}) for the running tree $T$ and
for \thmref{c09-dpjoin} -- the pair/peak decomposition proved once for
diameter and generalised here -- and Chapter~2 (\chref{c02}) for the
bottom-up recursion pattern whose general-tree upgrade this whole chapter is.
\end{chaptertheorybox}

\begin{figure}[H]
\centering
\begin{tikzpicture}[diagram, x=9mm, y=7mm]
  \foreach \st/\lab in {1/lv1,2/lv2,3/lv3,4/lv4,5/lv5,6/lv6} {
    \fill[\lab!22, draw=\lab, line width=0.8pt] (\st-1,0) rectangle (\st,\st*0.52);
    \node[font=\scriptsize\bfseries, text=\lab] at (\st-0.5,\st*0.52-0.42) {L\st};}
  \node[annot, anchor=south west, font=\scriptsize]
    at (0.05,0.62) {13.1 child-fold template; state per node};
  \node[annot, anchor=south west, font=\scriptsize]
    at (1.05,1.14) {13.1 LC 333 subtree-BST tuple};
  \node[annot, anchor=south west, font=\scriptsize]
    at (2.05,1.66) {13.2 in/out pair; CF 1529C endpoints};
  \node[annot, anchor=south west, font=\scriptsize, text=lv5]
    at (3.05,2.18) {13.3 prop:c10-pairtop; harvest arms};
  \node[annot, anchor=south west, font=\scriptsize, text=lv5]
    at (4.05,2.70) {13.4 CF 161D distance-k; ex:c10-k2 census};
  \node[annot, anchor=south west, font=\scriptsize, text=lv6]
    at (5.05,3.22) {summit (CF 1076E, depth-capped subtree)};
\end{tikzpicture}
\caption*{Chapter~13 on the six-step ladder: the child-fold and LC~333 sit
at $\lv2$--$\lv3$; the in/out pair and CF~1529C climb to $\lv3$--$\lv4$;
\propref{c10-pairtop} and the distance-$k$ histogram are $\lv4$ work; the
Summit (CF~1076E, subtree-of-$v$ depth-capped counting via a difference
array on the Euler view) is step~6 territory viewed from step~5. Cookbook
path $=$ steps 2--3; full path runs to 6.}
\end{figure}

% ---------------- The running example ----------------
\subsection*{The Running Example: One Ten-Node Tree, Four DP Angles}

\emph{One} rooted tree does all the work -- the same $T$ Chapter~12 measured,
now wearing a DP state. The nodes are $1,\dots,10$, the root is fixed at $1$
(a \emph{choice}, not a fact; rerooting is Chapter~15's promise, not this
chapter's), and the edges are
\[
  1\!-\!2,\; 1\!-\!3,\; 2\!-\!4,\; 2\!-\!5,\; 3\!-\!6,\; 3\!-\!7,\; 5\!-\!8,\; 7\!-\!9,\; 9\!-\!10 .
\]
Children of the root view: $1 \!\to\! \{2,3\}$, $2 \!\to\! \{4,5\}$,
$3 \!\to\! \{6,7\}$, $5 \!\to\! \{8\}$, $7 \!\to\! \{9\}$, $9 \!\to\! \{10\}$;
the leaves are $\{4,6,8,10\}$. Read the picture first (\figref{c10-runtree-dp}):
every badge on every node is a hand-computed entry of a table that gets reused
in \secref{c10-twostate} (in/out), \secref{c10-harvest} (peaks) and
\secref{c10-distance} (the $k=2$ census).

\begin{figure}[H]
\centering
\begin{tikzpicture}[diagram, scale=0.90, transform shape,
    tnode/.style={circle, draw=primaryblue, line width=0.9pt, fill=primaryblue!8,
                  minimum size=8mm, font=\small\bfseries, text=primaryblue!60!black}]
  \node[tnode] (n1) at (0,0)     {$1$};
  \node[tnode] (n2) at (-3,-1.6) {$2$};
  \node[tnode] (n3) at ( 3,-1.6) {$3$};
  \node[tnode] (n4) at (-5.2,-3.2) {$4$};
  \node[tnode] (n5) at (-1.6,-3.2) {$5$};
  \node[tnode] (n6) at ( 1.6,-3.2) {$6$};
  \node[tnode] (n7) at ( 5.2,-3.2) {$7$};
  \node[tnode] (n8) at (-1.6,-4.8) {$8$};
  \node[tnode] (n9) at ( 5.2,-4.8) {$9$};
  \node[tnode] (n10) at ( 5.2,-6.4) {$10$};
  \draw[tedge] (n1)--(n2) (n1)--(n3) (n2)--(n4) (n2)--(n5)
               (n3)--(n6) (n3)--(n7) (n5)--(n8) (n7)--(n9) (n9)--(n10);
  % in/out badges: dp_in above (successgreen), dp_out below (warningred)
  \foreach \n/\xin/\xout in {n1/6/5, n2/2/2, n3/2/3, n4/1/0, n5/1/0,
                             n6/1/0, n7/2/1, n8/1/0, n9/1/1, n10/1/0} {
    \node[font=\scriptsize\bfseries, text=successgreen!55!black, anchor=south]
      at (\n.north) {\scriptsize in=\xin};
    \node[font=\scriptsize\bfseries, text=warningred!70!black, anchor=north]
      at (\n.south) {\scriptsize out=\xout};
  }
  \node[annot, font=\scriptsize, text=difforange!70!black, anchor=west]
    at (7.4,0) {root $1$, fixed};
  \node[annot, font=\scriptsize, text=successgreen!55!black, anchor=west]
    at (7.4,-0.7) {$\mathit{dp}_{\mathit{in}}[v]$: best with $v$ taken};
  \node[annot, font=\scriptsize, text=warningred!70!black, anchor=west]
    at (7.4,-1.4) {$\mathit{dp}_{\mathit{out}}[v]$: best with $v$ skipped};
  \node[annot, font=\scriptsize, text=softpurple, anchor=west]
    at (7.4,-2.2) {MIS size $=\max(6,5)=6$};
  \node[annot, font=\scriptsize, text=softpurple, anchor=west]
    at (7.4,-2.85) {witness $S=\{1,4,6,7,8,10\}$};
\end{tikzpicture}
\caption{The running tree $T$ \emph{with its in/out DP overlay}. Rooted at
$1$, the two-state MIS recurrence (\thmref{c10-mis-tree}) computes
$\mathit{dp}_{\mathit{in}}[v] = 1 + \sum_{c}\mathit{dp}_{\mathit{out}}[c]$ and
$\mathit{dp}_{\mathit{out}}[v] = \sum_c \max(\mathit{dp}_{\mathit{in}}[c],\mathit{dp}_{\mathit{out}}[c])$
in a single post-order DFS -- the same table every section of this chapter
reuses. Peak census and distance-$2$ census on this picture live in
\secref{c10-harvest} and \exref{c10-k2}.}
\label{fig:c10-runtree-dp}
\end{figure}

Four angles, four hand-computed tables -- each lives in the section named on
its row and never changes:

\begin{center}
\small
\renewcommand{\arraystretch}{1.25}
\begin{tabular}{@{}p{0.28\linewidth} p{0.46\linewidth} l@{}}
\toprule
\textbf{Angle on $T$} & \textbf{What comes out} & \textbf{Traced in} \\
\midrule
child-fold of sizes & sum-of-subtree-size $= \sum_v sz(v) = 10\!+\!4\!+\!5\!+\!1\!+\!2\!+\!1\!+\!3\!+\!1\!+\!2\!+\!1 = 30$ & \secref{c10-subtrees} \\
in/out MIS & $\mathit{dp}_{\mathit{in}}=[6,2,2,1,1,1,2,1,1,1]$, $\mathit{dp}_{\mathit{out}}=[5,2,3,0,0,0,1,0,1,0]$, MIS $=6$ & \secref{c10-twostate} \\
peak census & every unordered pair $(u,v)$ harvests at $\ell=\mathrm{LCA}(u,v)$; $\binom{10}{2}=45$ pairs grouped by $\ell$ & \secref{c10-harvest} \\
distance-$k=2$ census & $10$ pairs at distance $2$; degenerate rows with one arm length~$0$: $(1,4),(1,5),(1,6),(1,7),(2,8),(3,9),(7,10)$ & \exref{c10-k2} \\
\bottomrule
\end{tabular}
\end{center}

% ---------------- Taxonomy chart ----------------
\begin{figure}[H]
\centering
\begin{tikzpicture}[diagram, scale=0.95, transform shape,
    mnode/.style={circle, draw=#1, line width=0.9pt, fill=#1!8, minimum size=6.5mm,
                  font=\scriptsize\bfseries, text=#1!60!black}]
  \node[mnode=primaryblue] (a) at (0,0)     {$\mathcal S(v)$};
  \node[annot, font=\scriptsize, anchor=west] at (0.9,0)     {$=$ one tuple per node, folded from children \quad (13.1)};
  \node[mnode=successgreen] (b) at (0,-1)   {$\mathcal I(v)$};
  \node[annot, font=\scriptsize, anchor=west] at (0.9,-1)   {$=$ two numbers, take/skip, combined by $\max$ \quad (13.2)};
  \node[mnode=difforange]   (c) at (0,-2)   {$\mathcal P(\ell)$};
  \node[annot, font=\scriptsize, anchor=west] at (0.9,-2)   {$=$ arms at the peak $\ell$, pair-count by \propref{c10-pairtop} \quad (13.3)};
  \node[mnode=warningred]   (d) at (0,-3)   {$H_v[\cdot]$};
  \node[annot, font=\scriptsize, anchor=west] at (0.9,-3)   {$=$ histogram over depths in $v$'s subtree, knapsack merge \quad (13.4)};
  \node[mnode=softpurple]   (e) at (0,-4)   {$\mathit{ans}$};
  \node[annot, font=\scriptsize, anchor=west] at (0.9,-4)   {$=$ read from the root (13.1--13.2) or sum over $\ell$ (13.3--13.4)$ $};
\end{tikzpicture}
\caption*{The DP taxonomy chart -- four shapes, one DFS. Every problem in
this chapter is one of these rows. The last two share the harvest-at-peak
identity \thmref{c09-dpjoin} (Ch.12) as their engine; the third is that
identity generalised, the fourth is it histogrammed.}
\end{figure}

% ---------------- Who owns what ----------------
\subsection*{Who Owns What (so nothing is learned twice)}

\begin{center}
\small
\renewcommand{\arraystretch}{1.25}
\begin{tabular}{@{}p{0.30\linewidth} p{0.32\linewidth} p{0.32\linewidth}@{}}
\toprule
\textbf{Idea} & \textbf{Owner} & \textbf{This chapter's share} \\
\midrule
height/size recursion on \emph{binary} trees & Ch.~1 (\thmref{c01-height-recursion}), Ch.~2 bottom-up & \emph{general} tree, tuple state, one DFS with a parent parameter \\
rooting an undirected tree & Ch.~4 (\chref{c04}) & assumed throughout; not re-taught \\
tree \emph{metrics} (depth, height, size, diameter, centroid, radius) & Ch.~12 (\chref{c09}) & treated as already-computed inputs to the DP \\
join-two-arms-at-the-peak for diameter & Ch.~12 (\thmref{c09-dpjoin}) & \propref{c10-pairtop} generalises to \emph{any} additive path functional \\
knapsack-over-children budget + sum-of-subtree-size bound & Ch.~14 (\chref{c27}) & 13.4 previews without proof \\
rerooting as a theory / ``answer for every root'' & Ch.~15 (\chref{c20}) & \emph{not} here -- root is fixed throughout \\
per-pair distance \emph{queries} via LCA & Ch.~16 (\chref{c11}) & one-shot DP; queries are what \secref{c10-distance} will make queryable there \\
Euler tour flattening (tin/tout) & Ch.~5 (\chref{c24}) & Summit CF~1076E uses the depth view, not a tour \\
DSU-on-tree & Ch.~17 (\chref{c28}) & not here -- this chapter's merges are per-child, not small-to-large \\
centroid decomposition & Ch.~24 (\chref{c18}) & not here -- Ch.~12 already closed the ``find one'' step \\
\bottomrule
\end{tabular}
\end{center}

% ---------------- Recognition guide ----------------
\begin{recognitionbox}[Pattern card -- ``the problem asks for a per-node (or per-pair) number computed from children'']
\reckey{the statement asks you to compute something \emph{at every node}
(the answer is $\mathcal S(v)$ for each $v$, or a $\sum_v$ over them), or
\emph{for every unordered pair} (a path functional $F(u,v)$ summed or counted
over all $\binom{n}{2}$ pairs).}
\recuse{
  \item ``compute for every node, one number'' (subtree size, subtree sum,
        subtree-max) -- the child-fold of \secref{c10-subtrees}, same
        machine as \thmref{c01-height-recursion} on a general tree;
  \item ``compute for every node, choose per node'' (take or skip, colour
        or leave, endpoint low or high) -- the in/out pair of
        \secref{c10-twostate}, CF~1529C's shape;
  \item ``count paths / pairs / sum over pairs satisfying $P(u,v)$'' -- the
        harvest-at-peak of \propref{c10-pairtop}, \secref{c10-harvest};
  \item ``count pairs at distance $k$ / sum path functionals with a length
        budget'' -- the depth-histogram of \secref{c10-distance}, CF~161D.
}
\recnotuse{
  \item the answer is for \emph{one root} only and no state travels up --
        that is Ch.~4's rooting idiom, not a DP cookbook;
  \item the tree is \emph{binary} and the state is constant (height, size,
        diameter) -- Chapters~1--2 already solve that; here the tree is
        general and the state is a \emph{tuple} or an \emph{array};
  \item the same question must be answered \emph{for every choice of root}
        (LC~834 family) -- that is rerooting, Chapter~15 (\chref{c20}), and
        this chapter will not solve it;
  \item you must answer many \emph{live pair queries} after preprocessing --
        hand it to LCA, Chapter~16 (\chref{c11}), and read \secref{c10-distance}
        first for the harvest identity they build on.
}
One-line test: \emph{does a fixed root turn the problem into ``each node
computes something from its children (or pairs of arms passing through
itself)''? Then it is a tree DP, and its cookbook is this chapter.}
\end{recognitionbox}

% ---------------- Constraint -> tool ----------------
\subsection*{Reading the Constraints (which DP shape the limits want)}

\begin{center}
\small
\renewcommand{\arraystretch}{1.25}
\begin{tabular}{@{}p{0.28\linewidth} p{0.40\linewidth} p{0.24\linewidth}@{}}
\toprule
\textbf{Constraint signature} & \textbf{What it rules in/out} & \textbf{Tool} \\
\midrule
$n \le 10^{5}$, tuple size $\bigO{1}$ per node & linear child-fold, no array on the state & \secref{c10-subtrees}, \secref{c10-twostate} \\
``assign a value in $[\ell_v, r_v]$ maximise $\sum |a_u - a_v|$'' & optimum at endpoints $\to$ two-state in/out & CF~1529C, \thmref{c10-mis-tree} \\
``count paths / pairs $(u,v)$ with property $P$'' & harvest at the peak over $\binom{n}{2}$ pairs & \propref{c10-pairtop} \\
$n \le 10^{5}$, $k \le 200$ (or $k$ small) & depth-histogram per node, $\bigO{nk}$ knapsack over children & \secref{c10-distance}, CF~161D \\
$n \le 10^{5}$, $k$ \emph{unbounded} & leave this chapter: the centroid-decomposition merge (Ch.~24, \chref{c18}) & \chref{c18} \\
``answer for every choice of root'' & rerooting, \emph{not} here & Ch.~15, \chref{c20} \\
$m \le 5\cdot10^{5}$ queries on subtrees with a depth cap & Summit shape: per-node DP plus difference array & CF~1076E \\
\bottomrule
\end{tabular}
\end{center}

% ---------------- Notation ----------------
\subsection*{Notation Used Only Here}

Root fixed at $1$ (or as stated); $c \in \mathit{ch}(v)$ ranges over the
children of $v$ in that rooted view. $\depth(v)$ is Ch.12's depth. $\mathcal
S(v)$ is the \emph{tuple} \secref{c10-subtrees} computes at $v$;
$(\mathit{dp}_{\mathit{in}}[v], \mathit{dp}_{\mathit{out}}[v])$ is the
two-state pair \secref{c10-twostate} keeps there;
$H_v[j] = |\{u \in \text{subtree}(v) : \depth(u) - \depth(v) = j\}|$ is the
depth histogram of \secref{c10-distance}. For a pair $(u,v)$, its
\term{peak} is $\ell = \mathrm{LCA}(u,v)$, and its two \term{arms} are
the path from $\ell$ down to $u$ and from $\ell$ down to $v$ (either may
be empty when $\ell$ equals $u$ or $v$). Distances count \emph{edges}. A
single-node tree has $\mathcal S$ computed at the leaf only, $\mathit{dp}_{\mathit{in}} = 1$,
$\mathit{dp}_{\mathit{out}} = 0$, and $H_v[0] = 1$.

Every section of this chapter is a chapter of the \emph{same} table on the
\emph{same} $T$: if a row here and a row there disagree, this opener is
wrong and the section is right -- fix it here, never the other way.
